\setcounter{secnumdepth}{3}
\titlespacing*{\chapter}{0pt}{-0.1cm}{1cm}
\titleformat{\chapter}[display]
  {\bfseries\LARGE}  
  {Phần~\thechapter}{0pt}{\raggedright}

\chapter{Phân tích heuristic (A*)}

\section{Tổng quan áp dụng A* cho bài toán Futoshiki}

Điểm quan trọng của việc áp dụng A* trong bài toán Futoshiki không nằm ở cấu trúc thuật toán, mà ở cách sử dụng heuristic để định hướng quá trình điền giá trị vào bảng.

Tại mỗi trạng thái, thuật toán thực hiện các bước sau:

\begin{itemize}
    \item \textbf{Chọn biến cần gán:} 
    Một ô chưa được điền sẽ được chọn để mở rộng. Việc lựa chọn này có thể dựa trên các tiêu chí như thứ tự duyệt hoặc ưu tiên các ô có nhiều ràng buộc.

    \item \textbf{Sinh các trạng thái con:}
    Với mỗi giá trị hợp lệ trong domain của ô đã chọn, tạo ra một trạng thái mới bằng cách gán giá trị đó vào ô tương ứng.

    \item \textbf{Áp dụng ràng buộc:}
    Sau khi gán, các ràng buộc được cập nhật (có thể kết hợp AC-3 hoặc forward checking) để thu hẹp domain của các ô còn lại.

    \item \textbf{Đánh giá heuristic:}
    Với mỗi trạng thái mới, giá trị $h(n)$ được tính toán dựa trên:
    \begin{itemize}
        \item Số lượng ràng buộc chưa thỏa mãn (Heuristic1), hoặc
        \item Mức độ chưa xác định của các biến thông qua domain (Heuristic2)
    \end{itemize}

    \item \textbf{Tính giá trị hàm đánh giá:}
    \[
    f(n) = g(n) + h(n)
    \]
    Trong đó:
    \begin{itemize}
        \item $g(n)$ là số bước đã gán (số ô đã điền thêm)
        \item $h(n)$ phản ánh mức độ còn lại cần hoàn thành
    \end{itemize}

    \item \textbf{Ưu tiên mở rộng:}
    Trạng thái có giá trị $f(n)$ nhỏ nhất sẽ được chọn để tiếp tục mở rộng.
\end{itemize}

Nhờ việc sử dụng heuristic, A* không duyệt tất cả các khả năng như brute-force mà tập trung vào các trạng thái “hứa hẹn” hơn, tức là các trạng thái:
\begin{itemize}
    \item Có ít vi phạm ràng buộc, hoặc
    \item Có domain đã được thu hẹp đáng kể
\end{itemize}

Điều này giúp giảm đáng kể số lượng trạng thái cần xét và tăng hiệu quả giải bài toán.

% ==================================================

\section{Heuristic 1 -- Đếm vi phạm ràng buộc}

\subsection{Ý tưởng}

Heuristic1 đánh giá trạng thái dựa trên số lượng ô vi phạm các ràng buộc bất đẳng thức.

Một ô được xem là vi phạm nếu tồn tại ít nhất một ràng buộc:
\begin{itemize}
    \item Không thỏa mãn bất đẳng thức, hoặc
    \item Chưa xác định do liên quan đến ô có giá trị $0$
\end{itemize}

\subsection{Định nghĩa}

\[
h(n) = \text{số ô vi phạm}
\]

\subsection{Mã giả}

\begin{tcolorbox}[colback=gray!5, colframe=black, title=Heuristic1]
\begin{verbatim}
FUNCTION Heuristic1(state, data):
    count ← 0

    FOR mỗi ô (i, j):
        IF tồn tại ràng buộc vi phạm:
            count ← count + 1

    RETURN count
\end{verbatim}
\end{tcolorbox}

\subsection{Chứng minh admissible}

Gọi:
\begin{itemize}
    \item $V$ là tập các ô bị vi phạm
    \item $|V|$ là số lượng ô vi phạm
\end{itemize}

Ta có:
\begin{itemize}
    \item Mỗi ô vi phạm cần được xử lý ít nhất một lần để thỏa mãn ràng buộc.
    \item Mỗi bước gán chỉ có thể thay đổi một ô.
\end{itemize}

Do đó, chi phí tối thiểu để đạt nghiệm thỏa mãn:

\[
h^*(n) \ge |V|
\]

Trong khi heuristic:

\[
h(n) = |V|
\]

Suy ra:

\[
h(n) \le h^*(n)
\]

\[
\Rightarrow \text{Heuristic1 là admissible}
\]


\subsection{Ví dụ minh họa}

\begin{tcolorbox}[colback=gray!5, colframe=black, title=Ví dụ 3x3]
\begin{verbatim}
Grid:
0   2   0
0   0   0
0   0   0

Constraints:
(1,1) < (1,2)
(1,2) > (2,2)
(2,1) > (3,1)
\end{verbatim}
\end{tcolorbox}

Các ô vi phạm: $(1,1), (1,2), (2,1), (2,2)$

\[
h(n) = 4
\]

\subsection{Nhận xét}

\begin{itemize}
    \item Ưu điểm: đơn giản, admissible.
    \item Nhược điểm: đánh giá còn thô, chưa phản ánh sâu cấu trúc bài toán.
\end{itemize}

% ==================================================

\section{Heuristic 2 -- AC-3 và miền giá trị}

\subsection{Ý tưởng}

Heuristic2 sử dụng thuật toán AC-3 để lan truyền ràng buộc và đánh giá mức độ chưa xác định của trạng thái.

\begin{itemize}
    \item Nếu domain rỗng $\rightarrow$ trạng thái vô nghiệm.
    \item Ngược lại, đếm số biến có nhiều hơn một giá trị khả dĩ.
\end{itemize}

\subsection{Định nghĩa}

\[
h(n) =
\begin{cases}
\infty & \text{nếu có domain rỗng} \\
\text{số biến có } |domain| > 1 & \text{ngược lại}
\end{cases}
\]

\subsection{Mã giả}

\begin{tcolorbox}[colback=gray!5, colframe=black, title=Heuristic2]
\begin{verbatim}
FUNCTION Heuristic2(state, data):
    domains ← copy(state.domains)

    IF AC3(domains) == FALSE:
        RETURN INFINITY

    count ← 0
    FOR mỗi domain d:
        IF size(d) > 1:
            count ← count + 1

    RETURN count
\end{verbatim}
\end{tcolorbox}

\subsection{Chứng minh admissible}

Heuristic2 được định nghĩa là số lượng biến có miền giá trị lớn hơn 1 sau khi áp dụng AC-3. Mặc dù trực quan cho thấy heuristic này phản ánh số lượng ô chưa được xác định, nhưng nó \textbf{không đảm bảo tính admissible}.
\subsubsection{Phản ví dụ}
Xét một trạng thái sau khi áp dụng AC-3:

\begin{center}
\begin{tabular}{|c|c|}
\hline
Biến & Domain \\
\hline
A & \{1,2\} \\
B & \{1,2\} \\
\hline
\end{tabular}
\end{center}

Giả sử tồn tại ràng buộc giữa $A$ và $B$ sao cho khi một biến được gán giá trị, biến còn lại sẽ được xác định duy nhất.

Khi đó:
\begin{itemize}
    \item Giá trị heuristic:
    \[
    h(n) = 2
    \]
    \item Tuy nhiên, chỉ cần một phép gán:
    \[
    A = 1 \Rightarrow B = 2
    \]
\end{itemize}

Suy ra chi phí thực tế:
\[
h^*(n) = 1
\]

Do đó:
\[
h(n) > h^*(n)
\Rightarrow \text{Heuristic2 không admissible}
\]

\subsubsection{Tính hữu ích của heuristic}

Mặc dù không admissible, Heuristic2 vẫn mang lại hiệu quả cao trong thực tế nhờ các đặc điểm sau:

\begin{itemize}
    \item \textbf{Phát hiện sớm trạng thái vô nghiệm:}  
    Khi AC-3 trả về domain rỗng, heuristic ngay lập tức loại bỏ trạng thái đó khỏi quá trình tìm kiếm.

    \item \textbf{Phản ánh mức độ chưa xác định:}  
    Số lượng biến có domain lớn hơn 1 cho thấy mức độ “chưa hoàn chỉnh” của trạng thái. Trạng thái có giá trị heuristic nhỏ thường gần nghiệm hơn.

    \item \textbf{Kết hợp với propagation:}  
    Việc sử dụng AC-3 giúp lan truyền ràng buộc trước khi đánh giá heuristic, từ đó giảm đáng kể không gian tìm kiếm.

    \item \textbf{Định hướng tìm kiếm hiệu quả:}  
    So với các heuristic đơn giản, Heuristic2 giúp A* ưu tiên các trạng thái có domain đã được thu hẹp, từ đó giảm số lượng trạng thái cần mở rộng.
\end{itemize}

\subsubsection{Kết luận}

Heuristic2 đánh đổi tính admissible để đạt được khả năng đánh giá trạng thái chính xác hơn trong thực tế. Mặc dù không đảm bảo tìm nghiệm tối ưu trong mọi trường hợp, heuristic này giúp cải thiện đáng kể hiệu năng tìm kiếm của thuật toán A*.
\subsection{Ví dụ minh họa}

Sau khi áp dụng AC-3 cho cùng ví dụ trên, ta có:

\begin{center}
\begin{tabular}{|c|c|c|}
\hline
Ô & Domain & Size \\
\hline
(1,1) & \{1\} & 1 \\
(1,2) & \{2\} & 1 \\
(1,3) & \{1,2,3\} & 3 \\
(2,1) & \{1,2,3\} & 3 \\
(2,2) & \{1\} & 1 \\
(2,3) & \{1,2,3\} & 3 \\
(3,1) & \{1,2,3\} & 3 \\
(3,2) & \{1,2,3\} & 3 \\
(3,3) & \{1,2,3\} & 3 \\
\hline
\end{tabular}
\end{center}

Có 6 biến có $|domain| > 1$

\[
h(n) = 6
\]

\subsection{Nhận xét}

\begin{itemize}
    \item Ưu điểm:
    \begin{itemize}
        \item Phát hiện sớm trạng thái vô nghiệm
        \item Phản ánh mức độ chưa xác định tốt hơn
    \end{itemize}
    \item Nhược điểm:
    \begin{itemize}
        \item Chi phí tính toán cao hơn
        \item Không đảm bảo admissible
    \end{itemize}
\end{itemize}

% ==================================================

\section{So sánh và đánh giá hai heuristic}

Để đánh giá hiệu quả của hai heuristic, ta xét khả năng \textbf{định hướng tìm kiếm} của chúng trong thuật toán A*.

\subsection{Hạn chế của Heuristic1}

Heuristic1 chỉ dựa trên việc đếm số ô vi phạm ràng buộc. Tuy nhiên, cách đánh giá này tồn tại một số hạn chế:

\begin{itemize}
    \item \textbf{Thông tin cục bộ:} Heuristic1 chỉ xem xét các ràng buộc bị vi phạm trực tiếp, mà không phản ánh được mức độ ràng buộc tổng thể của toàn bộ bảng.
    
    \item \textbf{Không phân biệt được các trạng thái tương đương:} Nhiều trạng thái có cùng số vi phạm nhưng mức độ ``gần nghiệm'' rất khác nhau.
    
\end{itemize}

\subsection{Ưu điểm của Heuristic2}

Heuristic2 sử dụng AC-3 để thu hẹp miền giá trị trước khi đánh giá trạng thái, do đó mang lại các ưu điểm sau:

\begin{itemize}
    \item \textbf{Khai thác thông tin toàn cục:} Thông qua AC-3, các ràng buộc được lan truyền trên toàn bộ bảng, giúp heuristic phản ánh chính xác hơn trạng thái hiện tại.
    
    \item \textbf{Phân biệt tốt hơn giữa các trạng thái:} Hai trạng thái có cùng số vi phạm có thể có domain rất khác nhau, và Heuristic2 có thể nhận ra sự khác biệt này.
    
    \item \textbf{Giảm không gian tìm kiếm:} Việc loại bỏ sớm các giá trị không hợp lệ giúp giảm số lượng trạng thái cần xét trong quá trình tìm kiếm.
\end{itemize}

\subsection{Ví dụ minh họa}

Xét hai trạng thái có cùng số vi phạm theo Heuristic1:

\begin{itemize}
    \item Trạng thái A: các domain đều lớn (nhiều giá trị khả dĩ)
    \item Trạng thái B: nhiều domain đã bị thu hẹp (gần xác định)
\end{itemize}

Khi đó:
\begin{itemize}
    \item Heuristic1: $h_1(A) = h_1(B)$ (không phân biệt được)
    \item Heuristic2: $h_2(A) > h_2(B)$ (phân biệt rõ trạng thái tốt hơn)
\end{itemize}

Do đó, A* sử dụng Heuristic2 sẽ ưu tiên trạng thái B, giúp tiến gần nghiệm nhanh hơn.

\subsection{Kết luận}

\begin{center}
\begin{tabular}{|c|c|c|}
\hline
Tiêu chí & Heuristic1 & Heuristic2 \\
\hline
Ý tưởng & Đếm vi phạm & Domain + AC-3 \\
Admissible & Có & Không đảm bảo \\
Độ chính xác & Thấp & Cao hơn \\
Chi phí tính & Thấp & Cao \\
\hline
\end{tabular}
\end{center}

Heuristic2 cung cấp thông tin giàu hơn so với Heuristic1 nhờ việc kết hợp constraint propagation. Mặc dù không đảm bảo tính admissible, heuristic này giúp định hướng tìm kiếm hiệu quả hơn, từ đó giảm số lượng trạng thái cần mở rộng và cải thiện thời gian giải bài toán trong thực tế.

Heuristic1 phù hợp khi cần đảm bảo tính tối ưu của A*, trong khi Heuristic2 giúp cải thiện hiệu năng tìm kiếm nhờ đánh giá trạng thái chính xác hơn.

Việc lựa chọn heuristic phụ thuộc vào mục tiêu giữa tối ưu nghiệm và thời gian thực thi.